\subsection{Esempio: lo spin di una particella in 3 dimensioni}
L'esempio più importante di spazio di Hilbert prodotto è dato dal grado di libertà di spin,
che, nella meccanica quantistica non relativistica, deve essere postulato (mentre emerge dalla
formulazione di Dirac).

$\vec{S}$ è un operatore vettoriale che segue l'algebra del momento angolare
$[S_i,S_j] = i\hbar\epsilon_{ijk}S_k$; dalla teoria del momento angolare si trova che, scelto il SCOC
$\{S^2,S_z\}$, gli autovalori degli autostati simultanei sono
$$
\begin{cases}
    S^2\ket{s,s_z} = \hbar^2s(s+1)\ket{s,s_z} & s = \frac{n}{2}, n\in\mathbb{Z} \\
    S_z\ket{s,s_z} = \hbar s_z\ket{s,s_z} & s_z = -s,-s+1,...,0,...,s-1,s\\
\end{cases}
$$
lo spazio di Hilbert dei gradi di libertà di spin è dunque uno spazio vettoriale di dimensione
$2s+1$ basato sul campo complesso: $\mathcal{H}_\text{spin} = \mathbb{C}^{2s+1}$.

Visto che questi operatori non "vedono" la parte spaziale, lo spazio di Hilbert totale di una
particella con spin $s$ in 3 dimensioni è
$$ \mathcal{H} = L^2(\mathbb{R}^3) \otimes \mathbb{C}^{2s+1}$$
Scelta come base $\{ \ket{\vec{x},s_z} = \ket{\vec{x}}\otimes\ket{s_z} \}$\footnote{con
$\ket{s_z = s} = e_1$ e $\ket{s_z=-s} = e_{2s+1}$ elementi della base standard di $\mathbb{C}^{2s+1}$},
si può scrivere il generico stato del sistema
$$\ket{\psi} = \sum_{s_z=-s}^{s}\int_{}^{}d^3\vec{x} \ket{\vec{x},s_z}\bra{\vec{x},s_z}\ket{\psi}
:= \sum_{s_z=-s}^{s}\int_{}^{}d^3\vec{x} \psi_{s_z}(\vec{x})\ket{\vec{x},s_z}$$
dunque
$$ \bra{\vec{x}}\ket{\psi}=\psi(\vec{x}) = \sum_{s_z=-s}^{s}\psi_{s_z}(\vec{x})\ket{s_z} =
    \begin{pmatrix}
	\psi_s(\vec{x}) \\
	\vdots \\
	\psi_{-s}(\vec{x})
    \end{pmatrix}$$
che è detta funzione d'onda vettoriale.

\section{Postulato 1 della Meccanica Quantistica}
\noindent\fbox{%
    \parbox{\textwidth}{%
    \textbf{I possibili risultati della misura di un osservabile F sono solo i suoi autovalori}
    $\{\mathbf{f_i}\}$\textbf{. Ciascun autovalore ha una probabilità}
    $\mathbf{\mathcal{P}_{\psi}(f_i)}$ \textbf{dipendente dallo stato del sistema.}
    }%
}
Questo è detto, per ovvii motivi, postulato della misura
